% =====================================================================
%  RESULTS SECTION -- prepared skeleton; numbers pending.
%  Every pending slot is marked \tbd (tables) or [figure placeholder].
% =====================================================================

% Lightweight placeholder box for figures not yet generated.
% Replace each \figbox{width}{height}{filename} with
%   \includegraphics[width=width]{filename}
% once the figure exists.
\newcommand{\figbox}[3]{%
  \setlength{\fboxsep}{2pt}%
  \fcolorbox{gray!55}{gray!12}{%
    \parbox[c][#2][c]{#1}{\centering\ttfamily\scriptsize
      [figure placeholder]\\[2pt]#3}}%
}

We evaluate the method on three cases of increasing realism: a linear--Gaussian system with closed-form posterior (no network trained, isolating the sampler from model error), stochastically forced 2D Navier--Stokes on a learned prior, and a multi-variable urban airflow. We instantiate the three samplers of Section~\ref{subsec:summary_table}, reported as \emph{Ours (SI-SDE/DM-SDE/FM-ODE)}.

\paragraph{Baselines.}
We compare against (i) generative posterior samplers sharing our trained prior -- FlowDAS \cite{chen_flowdas_2025}, SDA \cite{rozet_score-based_2023}, and D-Flow SGLD \cite{parikh_d-flow_2026,ben-hamu_d-flow_2024}; (ii) the SURGE sequential-Monte-Carlo particle filter \cite{wei_surge_2026}, which wraps a guided proposal with Girsanov-corrected importance reweighting, evaluated on top of the FlowDAS and SDA guidance; and (iii) conventional filters on the \emph{true} forward solver (Navier--Stokes only) -- the ensemble Kalman filter \cite{evensen_data_2022} and the bootstrap particle filter \cite{carrassi_data_2018}. Each method is described in Appendix~\ref{appendix:methods}.

\paragraph{Metrics and scenarios.}
We report, over an ensemble and averaged over steps and trajectories: \emph{(a) point accuracy} -- ensemble-mean RMSE and an energy-spectrum RMSE; \emph{(b) calibration} -- CRPS, spread--skill ratio, and rank histograms; \emph{(c) distributional fidelity} -- KL divergence to the reference marginal posteriors at observed and unobserved points (exact in the analytical case); \emph{(d) cost} -- network evaluations per step (NFE) and wall-clock, at matched ensemble size. Two linear observation operators structure the field-valued comparisons: \emph{super-resolution} ($32^2\!\to\!128^2$, $16^2\!\to\!128^2$ down-sampling) and \emph{sparse sensors} (random subsets at $5\%$ and $1.5625\%$), with additive Gaussian observation noise.

% =====================================================================
\subsection{Simple analytical test case}
\label{subsec:results_analytical}
% =====================================================================
\begin{wrapfigure}[11]{r}{0.46\linewidth}
    \centering
    \includegraphics[width=\linewidth]{figures/analytical/analytical_kl_vs_M.pdf}
    \vspace{-20pt}
    \caption{KL to the exact posterior vs.\ sampler steps $M$.}
    \label{fig:analytical}
\end{wrapfigure}

The first case is the linear--Gaussian system of Appendix~\ref{appendix:simple_test_case}: $\state^1 = \state^{0} + \modelnoise$, $\obs^1 = \obsmatrix\state^1 + \obsnoise$ with standard-normal noises and $\obsmatrix=\obscov=\bI$, so the exact Gaussian posterior is available and the SI drift is known in closed form (Prop.~B.9 of \cite{chen_probabilistic_2024}); no network is trained, isolating the sampler from model error.

The KL divergence to the exact posterior versus the number of sampler steps is shown in Figure~\ref{fig:analytical}; the prior, likelihood, exact and sampled posteriors, and KL convergence are in Figure~\ref{fig:analytical_panels} (Appendix~\ref{appendix:additional_results}).

% TODO(results): 2-3 sentences on the analytical findings once final numbers
% are in (inflated covariance vs. DPS-surrogate baselines; agreement across
% the three samplers; EnKF/PF as the linear-Gaussian reference).

% =====================================================================
\subsection{Stochastic Navier--Stokes}
\label{subsec:results_ns}
% =====================================================================
The second case is an incompressible fluid on the torus $\mathcal{T}^2=[0,2\pi]^2$, governed by the stochastically forced two-dimensional vorticity equation,
\begin{align}\label{eq:ns_equation}
    \rd\vorticity + \velocity\cdot\nabla\vorticity\,\rd t
    = \nu\Delta\vorticity\,\rd t - \alpha\vorticity\,\rd t + \varepsilon\,\rd\xi,
\end{align}
where $\velocity=\nabla^\perp\psi$ derives from the stream function $\psi$ solving $-\Delta\psi=\vorticity$, and $\rd\xi$ is temporally white forcing on selected Fourier modes. The stochastic forcing makes the transition density non-deterministic, so the prior $\conddist{\vorticity^{n}}{\vorticity^{n-1}}$ is a distribution the generative model must capture. We observe vorticity through $\obs=\obsoperator(\vorticity)+\obsnoise$, $\obsnoise\sim\mathcal{N}(0,0.05^2\bI)$.

Table~\ref{tab:ns_accuracy} reports accuracy, calibration, and cost across methods and all four scenarios; trajectory snapshots and diagnostics (RMSE over steps, energy spectra, rank histograms, and the divergence to the reference posterior -- a large-ensemble EnKF on the true solver) are in Figures~\ref{fig:ns_trajectories}--\ref{fig:ns_diagnostics} (Appendix~\ref{appendix:additional_results}).

% TODO(results): headline findings paragraph once the grid finishes
% (dense vs. sparse regimes; Jacobian-free vs. ensemble-shared inflated;
% sampler agreement; cost comparison).

\begin{table*}[ht]
    \centering
    \scriptsize
    \setlength{\tabcolsep}{3.5pt}
    \caption{Navier--Stokes: vorticity RMSE, CRPS, and spread--skill ($|1-\mathrm{spread}/\mathrm{skill}|$, $0=$ calibrated) across the four scenarios, with per-step cost (lower is better). Five held-out trajectories at $E=64$, $M=50$; ensemble-shared inflated covariance (Jacobian-free variant in Table~\ref{tab:ns_accuracy_jacfree}). The solver-free baselines share the prior; the conventional filters use the true solver. \tbd{} cells pending.}
    \label{tab:ns_accuracy}
    \begin{tabular}{l|cccc|cccc|cccc|c}
        \toprule
        & \multicolumn{4}{c|}{\textbf{Vorticity RMSE}}
        & \multicolumn{4}{c|}{\textbf{CRPS}}
        & \multicolumn{4}{c|}{\textbf{Spread--skill}}
        & \textbf{Cost} \\
        & $16^2$ & $32^2$ & $5\%$ & $\tfrac{1}{64}$ & $16^2$ & $32^2$ & $5\%$ & $\tfrac{1}{64}$ & $16^2$ & $32^2$ & $5\%$ & $\tfrac{1}{64}$ & s/step \\
        \midrule
        \multicolumn{14}{l}{\textit{Ours -- ensemble-shared inflated covariance}} \\
        Ours (SI-SDE) & \tbd & \tbd & \tbd & \tbd & \tbd & \tbd & \tbd & \tbd & \tbd & \tbd & \tbd & \tbd & \tbd \\
        Ours (DM-SDE) & \tbd & \tbd & \tbd & \tbd & \tbd & \tbd & \tbd & \tbd & \tbd & \tbd & \tbd & \tbd & \tbd \\
        Ours (FM-ODE) & \tbd & \tbd & \tbd & \tbd & \tbd & \tbd & \tbd & \tbd & \tbd & \tbd & \tbd & \tbd & \tbd \\
        \midrule
        FlowDAS            & \tbd & \tbd & \tbd & \tbd & \tbd & \tbd & \tbd & \tbd & \tbd & \tbd & \tbd & \tbd & \tbd \\
        FlowDAS + SURGE    & \tbd & \tbd & \tbd & \tbd & \tbd & \tbd & \tbd & \tbd & \tbd & \tbd & \tbd & \tbd & \tbd \\
        SDA                & \tbd & \tbd & \tbd & \tbd & \tbd & \tbd & \tbd & \tbd & \tbd & \tbd & \tbd & \tbd & \tbd \\
        SDA + SURGE        & \tbd & \tbd & \tbd & \tbd & \tbd & \tbd & \tbd & \tbd & \tbd & \tbd & \tbd & \tbd & \tbd \\
        D-Flow SGLD        & \tbd & \tbd & \tbd & \tbd & \tbd & \tbd & \tbd & \tbd & \tbd & \tbd & \tbd & \tbd & \tbd \\
        \midrule
        EnKF               & \tbd & \tbd & \tbd & \tbd & \tbd & \tbd & \tbd & \tbd & \tbd & \tbd & \tbd & \tbd & \tbd \\
        Particle filter    & \tbd & \tbd & \tbd & \tbd & \tbd & \tbd & \tbd & \tbd & \tbd & \tbd & \tbd & \tbd & \tbd \\
        \bottomrule
    \end{tabular}
\end{table*}

% =====================================================================
\subsection{Urban airflow over a building array}
\label{subsec:results_urban}
% =====================================================================
The third case is a multi-variable urban computational-fluid-dynamics flow (uDALES) over a building array: coupled velocity ($u,v,w$) and potential-temperature ($\theta_\ell$) fields on a domain with solid obstacles, whose wakes and plumes make the statistics strongly anisotropic; solid cells are excluded from all metrics and sensors. The case is \emph{generative-only} (no true forward solver, so no conventional filters) and \emph{sparse-only}; Table~\ref{tab:urban_accuracy} (Appendix~\ref{appendix:additional_results}) reports per-variable RMSE, CRPS, spread--skill, and cost. An ablation isolating the covariance, diffusion-strength, step-count, and ensemble-size axes, with trajectory figures for all three cases, is also in Appendix~\ref{appendix:additional_results}.

% TODO(results): 2-3 sentences on the urban findings once numbers are in.
